\section{Results}

% \begin{table*}[!ht]
%     \small
%     \centering
%     \begin{tabular}{|w{c}{0.05cm}|w{c}{0.3cm}||c|c||c|c|c||c|c||c|c|c|}
%     \hline
%         & & \multicolumn{5}{c||}{GPT-4o} & \multicolumn{5}{c|}{Llama}\\
%         \hline
%         & & \multicolumn{10}{c|}{$\Delta\textbf{rater}$} \\
%         \hline
%         & & \multicolumn{2}{c||}{$\textnormal{H}_{H} \to \textnormal{LLM}_H$}  & \multicolumn{3}{c||}{$\textnormal{H}_{A} \to \textnormal{LLM}_A$} & \multicolumn{2}{c||}{$\textnormal{H}_{H} \to \textnormal{LLM}_H$}  & \multicolumn{3}{c|}{$\textnormal{H}_{A} \to \textnormal{LLM}_A$}\\
%         \hline
%         P & C. & Full & 3ex & Sep. & Bat. & Edited& Full & 3ex & Sep. & Bat. & Edited \\
%         \hline
%         \multirow{6}{*}{1}& H. & $0.437^{\dagger}$ & $0.387$ & - & - & - & $0.576$ & $0.571$ & - & - & -\\
%         \cline{2-12}
%         & Ide. & -& -& $0.474^{\star}$ & $0.464$ & $0.552^{s,b\uparrow}$ & -& -& $0.552$ & $0.559$ & $0.552$\\
%         \cline{2-12}
%         & Or. & -& -& $0.497^{\star}$ & $0.448$ & $0.544^{s,b\uparrow}$ & -& -&$0.547^{\star}$ & $0.524$ & $0.529^{s\downarrow}$\\
%         \cline{2-12}
%         & WC & -& -& $0.452^{\star}$ & $0.439$ & $0.554^{s,b\uparrow}$ & -& -& $0.547^{\star}$ & $0.527$ & $0.507^{s\downarrow}$\\
%         \cline{2-12}
%         & SF & -& -& $0.459$ & $0.455$ & $0.553^{s,b\uparrow}$ & -& -& $0.559^{\star}$ & $0.512$ & $0.523^{s\downarrow}$\\
%         \cline{2-12}
%         & Cv. & -& -& $0.362$ & $0.388^{\star}$ & $0.472^{s,b\uparrow}$ & -& -& $0.491$& $0.477$ & $0.487$\\
%         \hline\hline
%         \multirow{5}{*}{4}& H. & $0.695^{\dagger}$ & $0.687$ & - & - & - & $0.700$ & $0.699$ & - & - & -\\
%         \cline{2-12}
%         & Ct. & -& -& $0.673$ & $0.696^{\star}$ & $0.698^{s\uparrow}$& -& -& $0.698^{\star}$ & $0.693$ & $0.701$\\
%         \cline{2-12}
%         & PA & -& -& $0.662$ & $0.677^{\star}$ & $0.680^{s\uparrow}$& -& -& $0.683^{\star}$ & $0.665$ & $0.682^{b\uparrow}$\\
%         \cline{2-12}
%         & La. & -& -& $0.627^{\star}$ & $0.595$ & $0.639^{s,b\uparrow}$& -& -& $0.630^{\star}$ & $0.573$ & $0.621^{b\uparrow}$\\
%         \cline{2-12}
%         & Na. & -& -& $0.669^{\star}$ & $0.652$ & $0.682^{s,b\uparrow}$& -& -& $0.674^{\star}$ & $0.642$ & $0.669^{b\uparrow}$\\
%         \hline\hline
%         \multirow{6}{*}{6} & H. & $0.629$ & $0.644^{\dagger}$ & - & - & -& $0.666$ & $0.680^{\dagger}$ & - & - & -\\
%         \cline{2-12}
%         & Ct. & -& -& $0.610$ & $0.652^{\star}$ & $0.676^{s,b\uparrow}$& -& -& $0.680^{\star}$ & $0.648$ & $0.694^{s,b\uparrow}$\\
%         \cline{2-12}
%         & PA & -& -& $0.605$ & $0.608$ & $0.666^{s,b\uparrow}$ & -& -& $0.652^{\star}$ & $0.601$ & $0.668^{s,b\uparrow}$\\
%         \cline{2-12}
%         & La. & -& -& $0.524^{\star}$ & $0.510$ & $0.546^{s,b\uparrow}$& -& -& $0.542^{\star}$ & $0.496$ & $0.547^{b\uparrow}$\\
%         \cline{2-12}
%         & Na. & -& -& $0.562$ & $0.557$ & $0.590^{s,b\uparrow}$& -& -& $0.579^{\star}$ & $0.547$ & $0.589^{b\uparrow}$\\
%         \hline\hline
%         & & \multicolumn{10}{c|}{$\Delta\textbf{rubric}$} \\
%         \hline
%         & & \multicolumn{2}{c||}{$\textnormal{H}_{H} \to \textnormal{H}_A$}& \multicolumn{3}{c||}{$\textnormal{LLM}_H \to \textnormal{LLM}_A$} & \multicolumn{2}{c||}{$\textnormal{H}_{H} \to \textnormal{H}_A$}& \multicolumn{3}{c|}{$\textnormal{LLM}_H \to \textnormal{LLM}_A$}\\
%         \hline
%         P & C. & Full & 3ex & Sep. & Bat. & Edited& Full & 3ex & Sep. & Bat. & Edited \\
%         \hline
%         \multirow{2}{*}{1} & full &$0.591$ &- & $0.756$& $0.750$& $0.839^{s,b\uparrow}$ &- &- & $0.685^{\star}$ & $0.671$& $0.711^{s,b\uparrow}$\\
%         \cline{2-12}
%         & 3ex &- &- & $0.792$& $0.789$& $0.838^{s,b\uparrow}$ &- &- & $0.736^{\star}$& $0.713$& $0.768^{s,b\uparrow}$\\
%         \hline
%         \multirow{2}{*}{4} & full &$0.656$ &- &$0.838^{\star}$ & $0.817$ & $0.876^{s,b\uparrow}$ &- &- & $0.854^{\star}$& $0.780$ & $0.869^{s,b\uparrow}$\\
%         \cline{2-12}
%         & 3ex&- &- & $0.847^{\star}$& $0.823$& $0.880^{s,b\uparrow}$ &- &- & $0.861^{\star}$& $0.784$& $0.877^{s,b\uparrow}$\\
%         \hline
%         \multirow{2}{*}{6} & full &$0.681$ &- & $0.781^{\star}$& $0.773$& $0.842^{s,b\uparrow}$ & -&- & $0.803^{\star}$& $0.748$& $0.864^{s,b\uparrow}$\\
%         \cline{2-12}
%         & 3ex&- &- & $0.804$& $0.801$& $0.857^{s,b\uparrow}$ &- &- & $0.813^{\star}$& $0.751$& $0.868^{s,b\uparrow}$\\
%         \hline
%         % \hline
%         % & & \multicolumn{10}{c|}{$\Delta\textbf{rater+rubric}$} \\
%         % \hline
%         % & & \multicolumn{2}{c||}{$\textnormal{LLM}_{H} \to \textnormal{H}_A$}& \multicolumn{3}{c||}{$\textnormal{H}_H \to \textnormal{LLM}_A$} & \multicolumn{2}{c||}{$\textnormal{LLM}_{H} \to \textnormal{H}_A$}& \multicolumn{3}{c|}{$\textnormal{H}_H \to \textnormal{LLM}_A$}\\
%         % \hline
%         % 1 & & $0.574^{\dagger}$& $0.539$& $0.279$ &$0.271$ & $0.401^{s,b\uparrow}$ & $0.601$& $0.605$& $0.467$ & $0.467$&$0.470$\\
%         % \hline
%         % 4& & $0.691$& $0.688$& $0.669$ &$0.668$ &$0.682^{s,b\uparrow}$& $0.674$& $0.674$& $0.710^{\star}$ & $0.674$ & $0.707^{b\uparrow}$\\
%         % \hline
%         % 6 & & $0.645$ & $0.653^{\dagger}$& $0.585$ &$0.595^{\star}$ & $0.644^{s,b\uparrow}$ & $0.660$& $0.670^{\dagger}$& $0.637^{\star}$ &$0.592$ & $0.671^{s,b\uparrow}$\\
%         % \hline
%     \end{tabular}
%     \caption{Kendall's $\tau$ results on AES with GPT-4o and Llama for $\Delta$rater. P. indicates the essay prompt, and C. indicates what ratings are being compared, with ideas, organization, word choice, sentence fluency, and conventions compared for prompt 1, and content, prompt adherence, language, and narrativity compared for prompts 4 and 6. $\tau$ is calculated with singular numerical values for $\Delta$rater and calculated through Pareto dominance comparison for preferences for $\Delta$rubric. Significance tests between separate (sep.), batch (bat.), and edited prompts are performed, where ${^s}$ and ${^b}$ in the edited prompt column represents significant differences with separate and batch prompts respectively. $\dagger$ is indicated next to comparisons that are significantly larger within holistic prompts. $\star$ is indicated next to comparisons that are significantly larger between separate and batch comparisons. $\uparrow$ and $\downarrow$ represent that the $\tau$ value for edited prompts is significantly larger or smaller respectively with the separate $s$ or batch $b$ prompts' $\tau$. The lack of any dagger, star, or arrow denotes no statistical significance. $\textnormal{H}$ is a shortened form for $\textnormal{Human}$. }
%     \label{tab:asap_results}
% \end{table*}


% \begin{table*}[h!]
%     \centering
%     \begin{tabular}{|c|w{c}{0.3cm}||c|c||c|c|c||c|c||c|c|c|}
%     \hline
%         & & \multicolumn{5}{c||}{GPT-4o} & \multicolumn{5}{c|}{Llama}\\
%         \hline
%         \multirow{3}{*}{\rotatebox[origin=c]{90}{$\Delta$rater}} & & \multicolumn{2}{c||}{Holistic} & \multicolumn{3}{c||}{Analytic} & \multicolumn{2}{c||}{Holistic} & \multicolumn{3}{c|}{Analytic}\\
%         \cline{3-12}
%          & & 0ex & 3ex & Sep. & Bat. & Edited & 0ex & 3ex& Sep. & Bat. & Edited\\
%         \cline{3-12}
%          & & $0.536$& $0.585$& $0.464^{\star}$& $0.167$&$0.471^{b\uparrow}$ &$0.470$& $0.578$& $0.445^{\star}$& $0.166$&$0.426^{b\uparrow}$\\
%         \hline\hline
%              \multirow{3}{*}{\rotatebox[origin=c]{90}{$\Delta$rubric}} & Ex. & \multicolumn{2}{c||}{$\textnormal{H}_{H} \to \textnormal{H}_A$}& \multicolumn{3}{c||}{$\textnormal{LLM}_H \to \textnormal{LLM}_A$} & \multicolumn{2}{c||}{$\textnormal{H}_{H} \to \textnormal{H}_A$}& \multicolumn{3}{c|}{$\textnormal{LLM}_H \to \textnormal{LLM}_A$}\\
%         \cline{2-12}
%         & 0ex & $0.534$&-& $0.640^{\star}$& $0.455$& $0.640^{b\uparrow}$ & $0.551$&- & $0.531^{\star}$& $0.285$& $0.551^{b\uparrow}$\\
%         \cline{2-12}
%          & 3ex &-&-& $0.623^{\star}$& $0.484$& $0.617 ^{b \uparrow}$ &- &- & $0.566^{\star}$& $0.326$& $0.623^{b\uparrow}$\\
%          % \hline
%         %  \multirow{2}{*}{\rotatebox[origin=c]{90}{$\Delta$r+r}}& & \multicolumn{2}{|c||}{$\textnormal{LLM}_H \to \textnormal{H}_A$}& \multicolumn{3}{c||}{$\textnormal{H}_H \to \textnormal{LLM}_A$}& \multicolumn{2}{c||}{$\textnormal{LLM}_H \to \textnormal{H}_A$}& \multicolumn{3}{c|}{$\textnormal{H}_H \to \textnormal{LLM}_A$}  \\
%         % \cline{3-12}
%         % & & $0.541$& $0.530$&  $0.474^{\star}$& $0.34$& $0.459 ^{b\uparrow}$ & $0.478$& $0.545$&  $0.421^{\star}$& $0.183$& $0.424^{b\uparrow}$\\
        
%         \hline
%     \end{tabular}
%     \caption{Kendall's $\tau$ results on IF with GPT-4o and Llama. $\Delta$rubric is calculated through instruction following ratio comparison for preferences. All other calculations and significance follow the methodology of Table~\ref{tab:asap_results}.}
%     \label{tab:if_results}
% \end{table*}

\subsection{Edited Rubrics}
% maybe i just want to group everything here

% edited > non-edited
When editing rubrics for autoraters to improve human-autorater agreement, it is important to provide examples and context and remove confirmation bias from analytic rubrics.

% When editing rubrics for autoraters to improve human-autorater agreement, it is important to provide more examples more similar to each data point for most task-specific (analytic IF) rubrics, while just providing examples that span the range of ratings is sufficient for increased agreement in general rubrics (most of AES, holistic IF).

% how to organize:
% 1. edited analytic rubrics improve correlation to humans (gpt AES) + self alignment (gpt AES, llama AES)
% gpt AES analytic (sig)
% llama AES analytic (sig)
% 2. separate matters more for llama, and IF than the examples we added
% llama AES analytic (most sig)
% llama IF analytic (all sig)
% gpt IF analytic (all sig)


% 3. on general rubrics, full is better, on task specific rubrics, edited is better
% gpt AES holistic
% llama AES holistic
% 4. adding examples when there are none helps general holistic rubrics
% gpt IF holistic
% llama IF holistic

% holistic 6 and IF are task specific

% gpt AES holistic
% - on general rubrics, full is better (sig); on task specific rubrics, edited is better (sig)

% gpt AES analytic
% - edited always better (sig)
% - always improves self alignment

% llama AES holistic
% - on general rubrics, full is better (no sig); on task specific rubrics, edited is better (sig)

% llama AES analytic
% - separate is good, examples don't add sig improvements (most of the time)
% - always improves self alignment

% gpt IF holistic
% - improvement no sig

% gpt IF analytic
% - separate is good, examples don't add sig improvements (higher)
% - separate is good, examples don't add sig improvements for self alignment either (lower)

% llama IF holistic
% - improvement no sig

% llama IF analytic
% - separate is good, examples don't add sig improvements (lower)
% - separate is good, examples don't add sig improvements for self alignment either (higher)

\textbf{Using GPT with edited analytic AES rubrics mostly significantly improves agreement with humans.} In Table~\ref{tab:asap_results} under $\Delta\textnormal{rater}$ under the ``edited" column, adding examples and context in the edited rubric always improved human-autorater agreement, significantly in the majority of cases, when using GPT, indicated by $s,b\uparrow$. These improvements range from 0.696 (batch prompt 4 content) to 0.698 (edited prompt 4 content) for the smaller, non-significant improvements, to 0.439 (batch prompt 1 word choice) to 0.554 (edited prompt 1 word choice), for the larger, significant improvements. While the analytic rubrics lacked explanations for the examples that the holistic AES rubrics provided, this statistically significant improvement suggests that the examples and context still provide necessary grounding for autorater scores.

\textbf{Edited analytic AES rubrics significantly improves both models' self-alignment but not with IF rubrics.} In addition to improving human-autorater alignment, adding examples and context in the edited rubric always significantly improves GPT and Llama's self-agreement with their scores on the original holistic rubric, although these two rubrics have inherent differences, seen in Table~\ref{tab:asap_results} under $\Delta\textnormal{rubric}$, indicated by $s,b\uparrow$. Interestingly, these agreements surpass human-human agreement with the two types of unedited rubrics (0.591, 0.656, 0.681 compared to 0.750, 0.817, 0.773 for GPT, and 0.671, 0.780, and 0.748 for Llama). This could be because the same autorater is used with both rubrics, whereas human ratings are from different individuals. This suggests that while the prompts superficially evaluate different criteria, the autorater converges to a consistent understanding of essay quality across both rubrics, which aligns with the overall rating objective, a convergence not observed to the same degree in human ratings.

However, adding examples of human analytic scores does not necessarily increase either autoraters' self-alignment on IF. Adding examples decreased $\tau$ with GPT-4o ($\tau$ from 0.640 to 0.640, and 0.623 to 0.617 between separate and edited rubrics in Table~\ref{tab:if_results}), but increased $\tau$ with Llama (0.531 to 0.551 and 0.566 to 0.623 respectively). This suggests that autoraters in this case do not achieve a unified understanding of instruction following across prompts. Despite this, the ratio aggregation method aligns more closely with the autorater's internal reasoning about output instruction adherence. With edited prompts, GPT-4o's alignment (but not Llama's) exceeded human alignment on both prompts, even though the human expert annotators across both rubrics remained consistent.

\textbf{Reducing confirmation bias in all analytic rubrics provides mostly significant improvement for both models.} 
Prior work has highlighted confirmation bias \citep{doi:10.1518/001872008X354183}, in which conceived judgment is reinforced as the task progresses. This bias could affect analytic rubrics, leading raters to assign low scores across all criteria if they initially believe the piece of text is of poor quality. \citet{lee-etal-2024-unleashing} has shown that using separate conversations to rate each sub-criteria tends to increase agreement over using a single conversation to rate all sub-criteria. However, this bias has not been statistically assessed for autoraters, and may affect autoraters differently than humans. 

Across the analytic rubrics on both tasks, the majority of comparisons show separate rubrics significantly outperforming the batched rubrics. With Llama on AES and both models on IF, edited rubrics do not consistently outperform separate rubrics. This suggests the separation of the individual criteria is more important than the examples for Llama. The scoring rubric itself may contribute to the inconsistency for GPT-4o between AES and IF. In AES, example essays answered the same writing prompt as the evaluated essay, and were accompanied by the same scoring rubric during evaluation. Conversely, in IF, the examples did not correspond to the same decomposed questions being scored. The rubrics in IF are task-specific and vary significantly across instructions. 

In addition, there is a trend of $\tau$ decreasing and then increasing based on criteria order in prompts 4 and 6, observed with both GPT-4o (0.696, 0.677, 0.595, 0.652 and 0.652, 0.608, 0.510, 0.557 respectively) and Llama (0.693, 0.665, 0.573, 0.642 and 0.648, 0.601, 0.496, 0.547 respectively). This suggests potential differences in bias between humans and autoraters. This may also suggest that autoraters struggle with rating word choice or language. Qualitative analysis of model outputs for instruction following reveal that batch ratings consistently demonstrated a bias with probability distributions heavily skewed towards ``yes" responses when answering yes/no decomposed questions. This contrasts with the separate API call approach, which more closely resembled  human response distributions, suggesting an underlying bias effect in batch prompts.

\textbf{Task-specific AES holistic rubrics show significant improvement when reducing examples, but general AES holistic rubrics do not.} Reducing the number of examples from the full set to three (3ex) tends to decrease $\tau$ for GPT on AES for the general holistic rubrics in essay prompts 1 and 4 ($\tau$ dropped from 0.437 to 0.387 and 0.695 to 0.687, respectively), and increase $\tau$ for the task-specific holistic rubric in essay prompt 6 ($\tau$ rose from 0.629 to 0.644). This could also be because the examples have explanations for essay prompts 1 and 4, but not for 6. The score explanations could be more influential than the scores alone when examples were originally provided, but this could also be due to the type of rubric.

\textbf{Adding examples for IF holistic rubrics show improvement.}
Although there is an increase in $\tau$ when adding examples to the holistic rubric in IF for both autoraters (0.536 to 0.585 and 0.470 to 0.578 for GPT and Llama respectively), it is not statistically significant.


%\jessica{Move any rater+rubric stuff to the appendix - However, when using Llama as the autorater for both AES and IF, and when using GPT-4o for IF, adding examples to analytic prompts does not always improve human-autorater alignment over both separate and batch prompts (Table~\ref{tab:if_results}, which shows both an increase in $\tau$ from 0.464 to 0.471 in $\Delta$ rater between the separate and edited rubrics, and a decrease in $\tau$ from 0.474 to 0.459 in $\Delta$ rater + rubric).}


% However, in Table~\ref{tab:tau_if_gpt_agree} for IF, the autorater had less alignment with humans two out of three times when given the edited prompt as compared to the separate condition with analytic prompts in $\Delta$ rater and both times in $\Delta$ rubric. In AES, the example essays provided all answered the same writing prompt as the essay being evaluated, and were provided with the same scoring rubric at evaluation time. However, in IF, the examples do not have the same decomposed questions as the output being scored. The rubrics here are task specific and vary greatly from instruction to instruction. 

% In $\Delta$ rubric, the ratio aggregation method agrees more with the internal model of how an autorater reasons about how well an output followed instructions. There is higher $\tau$ in all analytic rubric conditions (0.60, 0.36, 0.56) compared to just comparing human alignment when the rubric is varied (0.32), even though in this dataset the human experts who annotated the holistic and analytic rubric were the same. Even though examples were not provided to humans, it is important when examples were not provided to autoraters in the same format, the autoraters still had higher alignment with itself (0.36 compared to 0.32). Therefore, it is important to study aggregation methods for such analytic rubrics that align with human behavior.

% When editing rubrics to improve autorater and human alignment, it is important to either provide examples that are general enough to capture a variety of conditions if the prompt is task specific, or to provide general prompts with a multitude of examples and context.

\subsection{Decomposition Level}
% analytic > holistic

Analytic rubrics do not consistently outperform holistic rubrics in aligning autoraters with human judgments due to prompt complexity or aggregation methods. 

\textbf{Prompt complexity moderates effects. }In Table~\ref{tab:asap_results} for AES, under the $\Delta$rater condition for essay prompt 1, almost all analytic batch prompts had higher correlation with humans ($\tau$ of 0.464, 0.455, 0.448, 0.439, 0.388) than the full holistic prompt ($\tau$ of 0.437). Conversely in essay prompts 4 and 6, the full holistic prompt outperformed most analytic batch prompts ($\tau$ of 0.695 compared to 0.696, 0.677, 0.652, 0.595 and $\tau$ of 0.629 compared to 0.652, 0.608, 0.557, 0.510). This discrepancy may stem from the complexity of the holistic prompts. Essay prompt 1's holistic prompt is highly complex, involving multiple sub-criteria that contain complex decisions, whereas, essay prompts 4 and 6 have less complex prompts. Introducing analytic rubrics may increase evaluation complexity, which leads to lower $\tau$. 

\textbf{Aggregation methods influence agreement. }Another factor is the prompt's output. The Pareto dominance aggregation method, a conservative estimate of essay comparison, is highly sensitive to disagreements in any single sub-criteria. In IF, both GPT-4o and Llama performed worse with analytic prompts than with holistic prompts ($\tau$ of 0.167 to 0.536 and 0.166 to 0.470). This is surprising, given that IF analytic prompts are considered task-specific, which should provide more detailed information about the task, potentially leading to higher agreement. This may be due to the aggregation method used--the ratio of ``yes'' to ``no'' responses--which does not account for the varying weights of sub-criteria in the overall holistic evaluation. Therefore, holistic preferences, being more straightforward to calculate, may yield higher performance. Higher human-autorater agreement is not necessarily achieved by decomposing a holistic rubric into several analytic parts, rather it is more important to understand the complexity of each evaluation measure as well as the aggregation methods used.

\subsection{Agreement Level}
Both datasets are stratified by human agreement level since human inter-rater agreement may influence human-autorater agreement.

\paragraph{High human inter-rater agreement is important.} On IF holistic rubrics, human-autorater agreement with the consolidated score was significantly higher when all three human annotators agreed under all conditions, than if only two or no annotators agreed with each other. This held for IF analytic rubrics for both the separate and edited conditions, although the batch condition had an increase, it was not significant. For AES holistic rubrics, a significant increase was only observed for prompt 4, while a non-significant increase was observed for prompt 6. Detailed analyses can be found in Appendix~\ref{sec:agreement}.

% \subsubsection{Separate vs. Batch}
% Separating analytic criteria into distinct API calls could improve agreement between autorater and human ratings compared to sending all criteria in a single API call. Prior work has highlighted confirmation bias \citep{doi:10.1518/001872008X354183}, in which a conceived judgment is reinforced as the task progresses. This bias could affect analytic criteria, leading raters to assign low scores across all aspects of an essay if they initially perceive it as poor. However, this bias has not been quantitatively assessed for autoraters. This bias may affect autoraters in a different way than humans, and using separate calls should mitigate this bias and increase agreement with human ratings.

% \paragraph{Separate mostly performs better than Batch.} \citet{lee-etal-2024-unleashing} has shown that using separate conversations to rate each sub-criteria tends to increase agreement over using a single conversation to rate all sub-criteria. Separate calls significantly improve alignment with humans in the $\Delta$Rater condition for GPT-4o for half of the comparisons, specifically for three out of five sub-criteria in essay prompt 1 (0.474 from 0.464, 0.497 from 0.448, 0.452 from 0.439), two out of four sub-criteria in essay prompt 4 (0.627 from 0.595, 0.669 from 0.652), and one out of four sub-criteria in essay prompt 6 (0.524 from 0.510), shown in Table~\ref{tab:asap_results}. This increases to 10 out of 12 comparisons when using Llama. For autorater self-alignment in the $\Delta$rubric condition, separate calls consistently outperform batch calls, although not all results are statistically significant. Lastly, in the $\Delta$Rater + Rubric condition, separate calls perform significantly better in essay prompts 4 and 6 with Llama (0.710 from 0.674 and 0.637 from 0.592 respectively). In IF, separate calls consistently and significantly outperform batch calls across all three conditions.

% \paragraph{Positional Bias exists.} There is a trend of $\tau$ decreasing and then increasing based on criteria order in prompts 4 and 6, observed with both GPT-4o (0.696, 0.677, 0.595, 0.652 and 0.652, 0.608, 0.510, 0.557 respectively) and Llama (0.693, 0.665, 0.573, 0.642 and 0.648, 0.601, 0.496, 0.547 respectively). This suggests potential differences in position bias between humans and autoraters, or a more pronounced positional bias in  autoraters. This may also suggest that autoraters struggle with rating word choice or language. Qualitative analysis of model outputs for instruction following reveal that batch ratings consistently demonstrated a bias with probability distributions heavily skewed towards ``yes" responses when answering yes/no decomposed questions. This contrasts with the separate API call approach, which more closely resembled  human response distributions, suggesting an underlying bias effect in batch prompts.